% Source: https://artofproblemsolving.com/community/c3h2404526_zero_to_power_a_real_number
% Stable identity: aops-topic-2404526-q01
For a positive integer $a=n$, $0^n=0\cdot0\cdots0=0$. For $a=\dfrac pq$ with positive integers $p,q$, $0^{\dfrac pq}=\left(\sqrt[q]{0}\right)^p=0^p=0$, because the only real number whose $q$-th power is $0$ is $0$. For an irrational $a>0$, $x^a$ is defined by continuity, and $x^a\to0$ as $x\to0^+$ (for example, $0<x^a<x^{r}$ for $0<x<1$ and any rational $0<r<a$, and $x^r\to0$). So for every positive real $a$,
\[0^a=\boxed{0}.\]

If $a<0$, then $0^a$ would have to be $\dfrac{1}{0^{-a}}=\dfrac10$, which is undefined, so $0^a$ has no value. For $a=0$ the two natural patterns disagree: $x^0=1$ for every $x\ne0$, while $0^y=0$ for every $y>0$. For this reason $0^0$ is an indeterminate form for limits and is often left undefined, although in algebra and combinatorics it is usually defined to be $1$.
